% Exported from Math++. Compile with pdflatex.
\documentclass[11pt,a4paper]{article}
\usepackage[T1]{fontenc}
\usepackage[utf8]{inputenc}
\usepackage{lmodern}
\usepackage{amsmath,amssymb,amsthm,mathtools}
\usepackage{graphicx,booktabs,xcolor}
\usepackage[hidelinks]{hyperref}
\usepackage[margin=25mm]{geometry}
\theoremstyle{plain}
\newtheorem{theorem}{Theorem}
\newtheorem{lemma}[theorem]{Lemma}
\newtheorem{corollary}[theorem]{Corollary}
\newtheorem{proposition}[theorem]{Proposition}
\theoremstyle{definition}
\newtheorem{definition}[theorem]{Definition}
\newtheorem{example}[theorem]{Example}
\theoremstyle{remark}
\newtheorem*{remark}{Remark}
\setlength{\parindent}{0pt}\setlength{\parskip}{0.6em}
\begin{document}
\section*{Integrals}

Indefinite integrals get a constant of integration.

\begin{equation*}
  \int \sin^2 x \, dx = \frac{x}{2} - \frac{\sin\left(2 x\right)}{4} + C
\end{equation*}
\begin{equation*}
  \int x e^x \, dx = e^{x} x - e^{x} + C
\end{equation*}
\begin{equation*}
  \int \frac{1}{x^2 + 1} \, dx = \arctan\left(x\right) + C
\end{equation*}
Definite integrals give a number:

\begin{equation*}
  \int_0^1 x^2 \, dx = \frac{1}{3}
\end{equation*}
\begin{equation*}
  \int_0^{\pi} \sin x \, dx = 2
\end{equation*}
\begin{equation*}
  \int_1^e \ln x \, dx = 1
\end{equation*}
Ask for a decimal answer by writing \texttt{$\approx$ ?} instead of \texttt{= ?}:

\begin{equation*}
  \int_0^1 e^{-x^2} \, dx \approx 0.7468241328
\end{equation*}
\end{document}
